\documentclass[12pt, a4paper]{article}

%% ── Engine & fonts ──────────────────────────────────────────────────────────
\usepackage[T1]{fontenc}
\usepackage{mathpazo}
\usepackage{microtype}

%% ── Mathematics ─────────────────────────────────────────────────────────────
\usepackage{amsmath, amssymb, amsthm}
\usepackage{mathtools}

%% ── Layout ──────────────────────────────────────────────────────────────────
\usepackage[margin=1.25in]{geometry}
\usepackage{setspace}
\onehalfspacing
\setlength{\parindent}{1.5em}
\setlength{\parskip}{0pt}

%% ── Sectioning & headings ───────────────────────────────────────────────────
\usepackage{titlesec}
\titleformat{\section}{\large\bfseries}{\thesection.}{0.75em}{}
\titleformat{\subsection}{\normalsize\bfseries}{\thesubsection.}{0.75em}{}
\titleformat{\subsubsection}{\normalsize\itshape}{\thesubsubsection.}{0.75em}{}

%% ── Headers & footers ───────────────────────────────────────────────────────
\usepackage{fancyhdr}
\pagestyle{fancy}
\fancyhf{}
\fancyhead[L]{\small\itshape Recursive Containment and Deferred Closure}
\fancyhead[R]{\small\itshape Flyxion}
\fancyfoot[C]{\small\thepage}
\renewcommand{\headrulewidth}{0.4pt}

%% ── Theorem environments ────────────────────────────────────────────────────
\theoremstyle{plain}
\newtheorem{theorem}{Theorem}[section]
\newtheorem{proposition}[theorem]{Proposition}
\newtheorem{lemma}[theorem]{Lemma}
\newtheorem{corollary}[theorem]{Corollary}
\newtheorem{conjecture}[theorem]{Conjecture}

\theoremstyle{definition}
\newtheorem{definition}[theorem]{Definition}
\newtheorem{example}[theorem]{Example}

\theoremstyle{remark}
\newtheorem{remark}[theorem]{Remark}

%% ── Operators ───────────────────────────────────────────────────────────────
\DeclareMathOperator{\Open}{open}
\DeclareMathOperator{\Pop}{pop}
\DeclareMathOperator{\Meld}{meld}
\DeclareMathOperator{\Reframe}{reframe}
\DeclareMathOperator{\Reach}{Reach}

%% ── Misc ────────────────────────────────────────────────────────────────────
\usepackage{hyperref}
\hypersetup{colorlinks=true, linkcolor=black, citecolor=black, urlcolor=black}

%% ────────────────────────────────────────────────────────────────────────────
\begin{document}

\begin{titlepage}
\vspace*{2cm}
\begin{center}
{\LARGE\bfseries Recursive Containment and Deferred Closure}\\[1em]
{\large A Scope Dynamics Framework for Symbolic Cognition}\\[3em]
{\large Flyxion}\\[0.5em]
{\normalsize Independent Researcher}\\[3em]
{\normalsize\today}
\end{center}
\vfill
\begin{abstract}
\noindent
Symbolic cognition is typically modeled as sequential token processing over
static representational states. We argue that this picture is incomplete in
a fundamental way. The primitive structure underlying arithmetic evaluation,
recursive computation, narrative comprehension, hypnotic induction, harmonic
tension, and creative insight is not the token or the state, but the
\emph{semantic scope}: a bounded, potentially unresolved region of meaning
that persists until lawfully closed. We develop a formal framework, rooted in
the Spherepop calculus of nested containment, centered on a stack of open
semantic bubbles~$\Sigma_t = [B_1, \ldots, B_n]$, a load functional
$L(\Sigma)$, and a family of transformation operators: $\Open$, $\Pop$,
$\Meld_\pi$, $\Reframe^c_\phi$, and $\Reframe^e_\phi$. Each operator is
defined with explicit admissibility conditions; inadmissibility is shown to
correspond to recognizable cognitive and rhetorical pathologies. The framework
is stratified: a tree-structured containment relation suffices for the core
cases, while DAG containment is required for metaphor and cross-domain
analogy. We advance an operator completeness conjecture as the central structural claim,
examine the sheaf-theoretic reading of multi-parent gluing, and close with
the philosophical consequences for models of intelligence.
\end{abstract}
\end{titlepage}

%% ────────────────────────────────────────────────────────────────────────────
\tableofcontents
\newpage

%% ── §1 ──────────────────────────────────────────────────────────────────────
\section{The Hidden Geometry of Scope}

\subsection{The Problem of Linear Symbolic Models}

Contemporary theories of cognition, language, and computation share a
background assumption that is rarely made explicit: that meaning unfolds
sequentially, one symbol following another, each state superseding its
predecessor. Arithmetic is presented as a sequence of operations. Language
comprehension is modeled as left-to-right token integration. Generative
systems predict the next element conditioned on what has come before. Even
working memory research tends to frame cognitive load as the number of items
retained rather than the structure of dependencies among them.

This picture is not simply incomplete. It is geometrically misleading. The
structure that makes arithmetic evaluable, narratives coherent, and trance
possible is not sequential at all. It is hierarchical, containment-based, and
essentially topological. The purpose of this paper is to make that structure
explicit and to develop a formal framework adequate to its complexity.

\subsection{Hidden Parentheses and the Normalization Insight}

Consider the elementary expression
\[
  1 + 2 \times 3.
\]
A student learning order of operations is told: multiplication before
addition. But this rule, presented as an external convention, conceals a
deeper fact. The expression already has a containment structure. Written
honestly, it reads
\[
  \bigl(1 + (2 \times 3)\bigr).
\]
The multiplication scope is \emph{inside} the addition scope. Precedence is
not a rule imposed on a flat string. It is a consequence of spatial depth.
The inner scope must resolve before the outer scope can be evaluated, not
because of convention, but because the outer scope's resolution function has
the inner scope's output as an unresolved dependency.

This observation generalizes. Every well-formed arithmetic expression has a
latent containment structure. The Spherepop normalization procedure makes
that structure explicit. Given any expression, one proceeds as follows: first,
insert all parentheses implied by conventional precedence, making every scope
boundary visible; second, regard each parenthesized region as a closed
geometric object, a circle in the plane; third, lift these circles into
three-dimensional bubbles nested inside one another; fourth, evaluate from
the innermost bubble outward, where each evaluation event is an irreversible
collapse of a local scope into a result value that propagates upward to its
parent.

The resulting \emph{Spherepop Normal Form} of $1 + 2 \times 3$ is the
reduction sequence
\[
  \bigl(1 + (2 \times 3)\bigr)
  \;\longrightarrow\;
  (1 + 6)
  \;\longrightarrow\;
  (7)
  \;\longrightarrow\;
  7.
\]
Each arrow is an irreversible semantic event. Each event reduces the
unresolved load of the system by closing one scope. The final result is not
merely a number. It is the residue of a history of nested commitments.

\subsection{Scope as the Primitive of Symbolic Cognition}

The normalization procedure just described extends far beyond arithmetic.
Natural language contains nested scopes at every level of structure: relative
clauses embedded inside noun phrases, subordinate clauses inside main clauses,
direct speech inside narrative frames, hypotheticals inside arguments. The
sentence
\[
  \textit{The boy who carried the box that held the bird ran home}
\]
contains at least four nested scopes of reference, each contributing
unresolved expectations that must remain active in working memory until their
containing clause is syntactically closed. Comprehension is not reading a
flat string. It is navigating a dynamic containment structure.

Programming languages make this architecture visible in their syntax.
Lexical scoping, block structure, function calls, and exception handlers are
all mechanisms for creating, suspending, and closing semantic regions. The
call stack maintained by any executing program is a direct implementation of
the scope stack we will formalize below.

Our central claim, developed and defended throughout this paper, is the
following: \emph{scope} is the primitive structure of symbolic cognition.
Not tokens. Not propositions. Not states. The bounded, potentially unresolved
semantic region that persists until lawfully closed. This geometric
reinterpretation of symbolic structure is philosophically aligned with
approaches that privilege spatial and transformational intuition over purely
sequential symbolic manipulation, most notably in geometric approaches to
analysis that insist the underlying spatial structure of a domain is
primary and the algebraic surface notation derived \cite{Needham1997}.


%% ── §2 ──────────────────────────────────────────────────────────────────────
\section{The Formal Framework}

\subsection{Bubbles and Containment Structures}

\begin{definition}[Semantic Bubble]
A \emph{semantic bubble} is a tuple $B = (C, E, U)$ where $C$ is a
contextual binding set, $E$ is an expectation structure encoding what
resolutions are anticipated, and $U(B) \geq 0$ is a scalar measuring
unresolved semantic load. A bubble $B$ is \emph{open} when $U(B) > 0$ and
\emph{closed} when $U(B) = 0$.
\end{definition}

\begin{definition}[Containment Structure]
A \emph{containment structure} is a pair $\Sigma = (B, \prec)$ where $B$ is
a finite set of semantic bubbles and $\prec$ is a strict partial order on $B$
(irreflexive, transitive, antisymmetric). We write $B_i \prec B_j$ to mean
$B_i$ is contained within $B_j$, and $B_i \prec^* B_j$ for the transitive
closure. The \emph{covering relation} $B_i \lessdot B_j$ holds when $B_i
\prec B_j$ and there is no $B_k$ with $B_i \prec B_k \prec B_j$.
\end{definition}

In the foundational treatment of this paper (Sections~\ref{sec:operators}
through~\ref{sec:hypnosis}), we restrict to the case where $\prec$ induces a
rooted forest: every bubble has at most one immediate parent. This tree
condition is satisfied by arithmetic evaluation, recursive computation,
ordinary narrative nesting, and hypnotic induction. Section~\ref{sec:dag}
relaxes it to a DAG for the treatment of metaphor and creative
cross-domain import.

\subsection{The Scope Stack and Semantic Load}

When the containment structure is a rooted tree, the current active path from
root to frontier can be linearized as a stack.

\begin{definition}[Scope Stack]
The \emph{scope stack} at time $t$ is the sequence
\[
  \Sigma_t = [B_1, B_2, \ldots, B_n]
\]
where $B_1 \prec B_2 \prec \cdots \prec B_n$ is the active containment path
and $B_n$ is the currently attended scope. Each $B_i$ is open: $U(B_i) > 0$.
\end{definition}

\begin{definition}[Semantic Load Functional]
The \emph{semantic load} of a scope stack $\Sigma = [B_1, \ldots, B_n]$ is
\[
  L(\Sigma) = \sum_{i=1}^{n} w_i \, U(B_i)
\]
where $w_i = f(d_i, s_i, r_i, c_i)$ is a weight combining stack depth $d_i$,
emotional salience $s_i$, recency $r_i$, and causal centrality $c_i$.
\end{definition}

The decomposition of $w_i$ into four partially independent dimensions is
essential. Semantic dominance is not purely hierarchical. A shallow but
emotionally charged bubble may exert greater gravitational pull on attention
than a deeper but neutral one. The weight function $f$ is left
underspecified here; what matters structurally is that the dimensions can
dissociate, producing the observed phenomena of shallow trauma overriding
deep abstraction, emotionally salient interruptions disrupting logical flow,
and central narrative frames organizing subordinate bubbles.

\subsection{Admissibility and the Resolution Rule}

The resolution rule is not a component of the bubble tuple but a
transformation operator defined separately. This objects-versus-morphisms
separation is deliberate: bubbles are semantic regions; operators are lawful
transformations over semantic regions.

\begin{definition}[Resolution Operator]
For each bubble $B \in \Sigma$, the \emph{resolution operator} is a partial
function
\[
  \rho : B \times \Sigma \rightharpoonup B'
\]
where $\rho(B, \Sigma)$ is defined if and only if all bubbles $B_j$ with
$B_j \prec B$ (all interior descendants of $B$ in the containment order)
satisfy $U(B_j) = 0$. Note that $B'$ denotes the updated bubble after
resolution, not a descendant; we use primed notation throughout for
post-operation states.
\end{definition}

The partiality is not a technical convenience. It is the formal expression
of the framework's central invariant.

\begin{definition}[Well-Nestedness]
A scope stack $\Sigma$ is \emph{well-nested} with respect to a candidate
pop operation at position $i$ if all $B_j$ with $B_j \prec B_i$ satisfy
$U(B_j) = 0$.
\end{definition}

\begin{proposition}
Well-nestedness is a necessary condition for the admissibility of any
resolution operation. A parent scope cannot lawfully close while any
descendant scope remains open.
\end{proposition}

This proposition, whose proof follows directly from the domain condition of
$\rho$, unifies arithmetic evaluation (innermost subexpression first),
recursive function return (callee before caller), proof completion (lemma
before theorem), narrative closure (embedded tale before framing tale), and
hypnotic release (innermost story before outer story). In every case, a
parent cannot lawfully resolve while dependent subordinates remain open.

The listener's phenomenological sense of \emph{unfinishedness} is
mathematically meaningful under this framework: the bubble literally lacks a
defined collapse morphism.

\subsection{Persistence and Suspended Activity}

A semantic bubble does not cease to exert causal influence merely because it
is no longer the foreground scope. Suspended scopes remain dynamically active
through their contribution to the semantic load functional $L(\Sigma)$ and
through unresolved dependency relations propagating upward through the
containment structure. This persistence distinguishes semantic bubbles from
ordinary stack frames in classical computational models. A conventional
execution frame becomes causally inert once control leaves it, except insofar
as it preserves local state for eventual return. In symbolic cognition, by
contrast, suspended scopes continue exerting attentional and affective
pressure even while subordinate scopes temporarily dominate conscious focus.
The unresolved narrative introduced at the opening of a story continues to
shape expectation while intervening events unfold. A suspended emotional frame
continues organizing perception even when not explicitly attended. An open
proof obligation continues constraining mathematical reasoning while auxiliary
lemmas are explored. Persistence is therefore not accidental to the framework
but constitutive of it. Containment structures are properly interpreted not as
static trees or DAGs but as temporally evolving semantic fields whose
unresolved regions remain causally active across time until lawfully collapsed.
This interpretation is compatible with dynamical accounts of cognition in which
causation operates through evolving constraint structures rather than through
isolated efficient interactions \cite{Juarrero1999}. More broadly, this connects
the present formalism to event-history ontologies in which the world is
constituted not by static states but by accumulated irreversible commitments
whose consequences propagate forward through time.


%% ── §3 ──────────────────────────────────────────────────────────────────────
\section{The Operator Calculus}
\label{sec:operators}

\subsection{Five Primitive Operators}

We define five operators on containment structures, each with explicit
admissibility conditions.

\begin{definition}[Open]
\[
  \Open(B) : \Sigma \mapsto \Sigma \cdot B
\]
Opens a new scope by appending $B$ to the current stack. $B$ must be a fresh
bubble not already in $\Sigma$. $U(B) > 0$ (newly opened scopes are by
definition unresolved).
\end{definition}

\begin{definition}[Pop]
\[
  \Pop(B_n) : [B_1, \ldots, B_{n-1}, B_n] \mapsto [B_1, \ldots, B_{n-1}']
\]
Closes the topmost scope. Admissible only when $\Sigma$ is well-nested at
position $n$. The parent $B_{n-1}$ is modified to $B_{n-1}'$, reflecting the
propagation of the resolved content upward. Closure is not isolated: it
changes the parent.
\end{definition}

\begin{definition}[Meld]
\[
  \Meld_\pi(B_i, B_j) : \Sigma \mapsto \Sigma'
\]
where $B_i$ and $B_j$ are sibling scopes (sharing a common parent but neither
containing the other), $\pi$ is a compatibility policy, and $B_i \sim_\pi
B_j$ denotes compatibility under $\pi$. The result $\Sigma'$ replaces $B_i$
and $B_j$ with a single merged bubble $B_{ij}$ at the same depth. Admissible
only when $B_i \sim_\pi B_j$; inadmissible when compatibility fails.
\end{definition}

Pop and Meld differ in a fundamental way. Pop is strictly causal and
sequential: it closes what was opened, in order. Meld is convergent and
cross-scope: it identifies sibling regions through shared resolution
structure. A joke callback is a Meld, not a Pop. The punchline does not
merely close the most recent frame; it retroactively identifies an earlier
frame with the present one, producing simultaneous collapse across distant
narrative regions. Musical cadence often behaves the same way, retroactively
gathering several unresolved harmonic scopes into a unified resolution rather
than closing them one by one.

\begin{definition}[Conservative Reframe]
\[
  \Reframe^c_\phi : (B, \prec) \to (B, \prec')
\]
where $\phi$ acts as the identity on $B$, $\prec'$ is a strict partial order,
and $\prec^* = \prec'^*$. The transitive closure is preserved exactly;
only the covering relation changes. Conservative reframe redistributes
containment without creating or destroying ancestral dependencies.
\end{definition}

\begin{definition}[Expansive Reframe]
\[
  \Reframe^e_\phi : (B, \prec) \to (B, \prec')
\]
where $\phi$ acts as the identity on $B$, $\prec'$ is a strict partial order,
and $\prec^* \subseteq \prec'^*$. The transitive closure is extended; new
ancestral dependencies may be introduced. Expansive reframe imports new
containment relations from outside the existing structure.
\end{definition}

\begin{definition}[Admissibility of Reframe]
A reframe $\Reframe_\phi$ of either species is \emph{admissible} only if
$\prec'$ is acyclic: for no $B_i$ does $B_i \prec'^* B_i$.
\end{definition}

Conservative reframe captures insight, therapeutic reinterpretation, and plot
twists: the same historical events become differently organized in the
containment hierarchy without any new facts being introduced. Expansive
reframe captures metaphor, analogy, and scientific paradigm shift: a semantic
region begins participating in a containment structure that was previously
foreign to it.

\subsection{Inadmissibility and Cognitive Pathology}

The admissibility conditions are not merely technical safeguards. Each
failure mode corresponds to a recognizable cognitive or rhetorical pathology.

\begin{proposition}[Inadmissibility Taxonomy]
The following four conditions are inadmissible and correspond to the named
phenomena.

(i) $\Pop(B_i)$ applied when some $B_j \prec B_i$ has $U(B_j) > 0$:
\emph{premature closure}. The narrative resolves before its embedded
obligations are discharged.

(ii) $\Meld_\pi(B_i, B_j)$ applied when $B_i \not\sim_\pi B_j$:
\emph{forced synthesis}. Two incompatible frames are collapsed as though
compatible, producing spurious apparent coherence.

(iii) $\Reframe_\phi$ producing a cycle $B_i \prec'^* B_i$:
\emph{semantic self-containment}. A scope depends on itself, producing
rumination, paradox, or incoherent circular explanation.

(iv) $\Reframe^e_\phi$ introducing ancestors incompatible with existing
gluing conditions across parent scopes: \emph{incoherent metaphor}. The
new containment relation cannot be made simultaneously consistent with all
existing parent structures.
\end{proposition}

There is also a dual failure for conservative reframe worth naming
separately. A $\Reframe^c_\phi$ that preserves the transitive closure but
severs all covering relations produces a structure in which every bubble is
ancestrally related to every other but nothing is locally adjacent to
anything. This is \emph{semantic dissociation}: global coherence with no
local navigability.

\subsection{Two Fundamental Theorems}

\begin{theorem}[Termination of Admissible Pop Sequences]
Every admissible Pop sequence over a finite well-nested tree containment
structure terminates in finitely many steps.
\end{theorem}

\begin{proof}
Define the \emph{depth} of a containment structure $(B, \prec)$ as
\[
  \delta(\Sigma) = \max \{ k : \exists\, B_{i_1} \prec B_{i_2} \prec \cdots \prec B_{i_k} \}.
\]
Since $B$ is finite, $\delta(\Sigma)$ is a non-negative integer. Each admissible
Pop operation removes the topmost bubble $B_n$ from the active stack and reduces
its parent's unresolved load, strictly decreasing $|B|$ by one. Because $|B|$
is finite and strictly decreasing under every admissible Pop, and because the
well-nestedness condition ensures no Pop can be applied to a bubble with
unresolved descendants (preventing indefinite suspension), the sequence of
admissible Pops must reach the empty stack $\Sigma = [\,]$ in at most $|B|$
steps.
\end{proof}

\begin{theorem}[Conservative Reframe Preserves Acyclicity]
If $(B, \prec)$ is a strict partial order and $\Reframe^c_\phi$ is an admissible
conservative reframe, then $(B, \prec')$ is also a strict partial order.
\end{theorem}

\begin{proof}
We verify that $\prec'$ satisfies irreflexivity, asymmetry, and transitivity.

\emph{Irreflexivity.} Suppose for contradiction that $B_i \prec' B_i$ for
some $B_i$. Then $B_i \prec'^* B_i$, which contradicts the admissibility
condition that $\prec'$ be acyclic.

\emph{Transitivity.} If $B_i \prec' B_j$ and $B_j \prec' B_k$, then
$B_i \prec'^* B_k$ by definition of transitive closure. Since $\prec^* =
\prec'^*$ (the conservative reframe condition), we have $B_i \prec^* B_k$.
The structure of $\prec'$ is chosen to realize this ancestry through some
covering path, so $B_i \prec' B_k$ or $B_i \prec'^* B_k$ through
intermediate nodes; either way transitivity holds.

\emph{Asymmetry.} If both $B_i \prec' B_j$ and $B_j \prec' B_i$, then
$B_i \prec'^* B_i$, contradicting acyclicity. Hence $\prec'$ is asymmetric.

Since $\prec'$ is irreflexive, asymmetric, and transitive, it is a strict
partial order.
\end{proof}

\begin{theorem}[Monotonicity of Semantic Load Under Pure Pop Sequences]
Let $\Sigma_0 \to \Sigma_1 \to \cdots \to \Sigma_n$ be a sequence consisting
only of admissible Pop operations with non-negative weights $w_i \geq 0$.
Then
\[
  L(\Sigma_{k+1}) \leq L(\Sigma_k)
\]
for all $k \in \{0, \ldots, n-1\}$.
\end{theorem}

\begin{proof}
Each admissible Pop removes the topmost open bubble $B_i$ satisfying
$U(B_i) > 0$, contributing $w_i U(B_i) > 0$ to $L(\Sigma_k)$. Propagation
effects modify the parent bubble $B_{i-1}$ to $B_{i-1}'$, but by the
admissibility condition all descendants of $B_i$ are already closed
($U(B_j) = 0$ for $B_j \prec B_i$), so no new unresolved load is introduced
by the propagation. Since $w_i \geq 0$ and $U(B_i) > 0$, removing $B_i$
strictly reduces the sum $\sum_j w_j U(B_j)$, yielding $L(\Sigma_{k+1}) <
L(\Sigma_k)$. In the degenerate case $U(B_i) = 0$ the bubble is already
closed and the Pop is trivially admissible with no change to $L$, giving
$L(\Sigma_{k+1}) = L(\Sigma_k)$. In both cases $L(\Sigma_{k+1}) \leq
L(\Sigma_k)$.
\end{proof}

\begin{remark}
Monotonicity fails for sequences containing Meld or Reframe operations.
Meld may temporarily increase local load during cross-scope identification
before convergence reduces it. Expansive Reframe by definition introduces
new ancestral dependencies, potentially increasing $L(\Sigma)$ before the
new structure stabilizes. This asymmetry reflects the phenomenology of
creative work and therapeutic restructuring: the path to lower total load
sometimes requires passing through higher-load intermediate configurations.
\end{remark}

\subsection{Operator Completeness}

\begin{conjecture}[Operator Completeness]
Every admissible transformation of a containment structure $(B, \prec)$ that
preserves node identity can be expressed as a finite composition of operations
drawn from $\{\Open, \Pop, \Meld_\pi, \Reframe^c_\phi, \Reframe^e_\phi\}$.
\end{conjecture}

The qualifier \emph{preserves node identity} is essential: it excludes
deletion (forgetting) and fabrication (confabulation), which are genuine
cognitive operations but lie outside the class of lawful semantic
transformations considered here. The boundary of this conjecture identifies
precisely where memory failure and motivated reasoning begin.

\subsection{Collapse Quotients and Irreversibility}

The Pop operator is irreversible in a structural rather than merely temporal
sense. When a scope collapses, its internal generative history is compressed
into a representative semantic residue. The resulting structure preserves
resolution outcome while discarding recoverable intermediate organization.
This compression can be formalized through a quotient construction. Let
$\sim_\rho$ be the equivalence relation on bubble histories induced by
admissible resolution under collapse operator $\rho$: two histories are
equivalent when they produce the same admissible semantic representative. The
collapsed bubble is then the quotient $B/{\sim_\rho}$.

The expression $(2 + 3)$ and the value $5$ illustrate the distinction
concretely. The unresolved scope $(2 + 3)$ and the collapsed residue $5$ are
extensionally equivalent but not structurally equivalent. One preserves
generative history; the other preserves only the lawful result of that
history. Abstraction is therefore not merely symbolic representation. It is
the quotienting of semantic histories by admissible collapse relations. This
account of abstraction as quotient connects the formal framework to the
broader theme of irreversibility that runs through the paper: what is lost in
a collapse is not content but provenance. The bubble $5$ is authoritative
precisely because it no longer carries the contingency of the history that
produced it. It excludes
deletion (forgetting) and fabrication (confabulation), which are genuine
cognitive operations but outside the scope of lawful semantic transformation
as defined here. The boundary of this conjecture thus identifies precisely
where memory failure and motivated reasoning begin.


%% ── §4 ──────────────────────────────────────────────────────────────────────
\section{Conversational Hypnosis as Scope Engineering}
\label{sec:hypnosis}

\subsection{The Nested Narrative Induction Pattern}

Conversational hypnosis, as practiced and systematized in the tradition
descended from Milton Erickson and formalized in later pedagogical work on
indirect induction, relies centrally on a structural technique: the
deliberate creation and suspension of nested narrative scopes. A practitioner
opens a story, interrupts it before closure to open a second story, interrupts
that to open a third, and then resolves them in reverse order. The listener's
experience is one of deepening absorption, suspended between multiple
unfinished narrative worlds.

What the Spherepop framework reveals is that this technique is not a
rhetorical trick layered on top of cognition. It is a direct exploitation of
cognition's native scope-management architecture.

\subsection{Formal Model of Induction}

Let the listener's active meaning-state at time $t$ be the scope stack
\[
  \Sigma_t = [B_1, B_2, \ldots, B_n]
\]
where each $B_i$ is an open semantic bubble: an unfinished story, image,
question, expectation, or emotional frame. Opening a narrative is a push:
\[
  \Open(B) : \Sigma \mapsto \Sigma \cdot B.
\]
Interruption before closure does not delete the suspended story. It suspends
it below the new scope:
\[
  [B_1] \xrightarrow{\Open(B_2)} [B_1, B_2].
\]
After several interruptions the listener carries
\[
  \Sigma = [B_1, B_2, B_3, \ldots, B_n].
\]
Resolution proceeds outward under the well-nestedness constraint, popping the
top scope first:
\[
  [B_1, B_2, B_3]
  \;\longrightarrow\;
  [B_1, B_2']
  \;\longrightarrow\;
  [B_1']
  \;\longrightarrow\;
  [\;].
\]
Each arrow is an irreversible semantic event. The final empty stack
corresponds to the completion of induction. The hypnotic release is the
outward cascade of scope completion.

\subsection{Trance as a Region in $\Sigma$-Space}

The semantic load functional $L(\Sigma) = \sum_i w_i U(B_i)$ provides a
precise characterization of the inductive state. Trance is not an altered or
mysterious state of consciousness requiring separate explanation. It is a
dynamically stable region in $\Sigma$-space characterized by four conditions:
high total unresolved semantic load $L(\Sigma)$, lawful nesting throughout,
deferred closure at every level, and maintained global coherence.

The fourth condition is critical and distinguishes trance from confusion. A
randomly incoherent stack would produce fragmentation. The hypnotic state
depends on the bubbles remaining globally navigable despite their suspension.
This is the sheaf condition operating on a tree: local sections (individual
bubble contents) must be mutually compatible even while global resolution is
deferred.

\subsection{Failed Induction as Failed Gluing}

Conversely, failed hypnotic induction corresponds formally to a failure of
the gluing condition. When the narrative scopes the practitioner opens cease
to overlap coherently, the listener cannot maintain global semantic
navigability. The subject drops out not because they resist, but because the
semantic manifold has fractured. The charts no longer paste together. At that
point no further deferred closure is possible because the structure is no
longer coherent enough to sustain it.

This analysis has a testable implication: the probability of successful
induction should be related to the practitioner's ability to maintain
pairwise compatibility between suspended narrative frames throughout the
induction sequence. Compatibility failure at any intermediate step propagates
upward and undermines the global structure.


%% ── §5 ──────────────────────────────────────────────────────────────────────
\section{Cross-Domain Instantiations}

\subsection{The Interpretation Map}

The formal structure developed in Section~\ref{sec:operators} is deliberately
domain-neutral. The scope stack $\Sigma_t$, the bubble tuple $B = (C, E, U)$,
and the operator family all admit interpretation in any domain where nested
unresolved regions are created, sustained, and eventually closed. What varies
across domains is the semantic realization of each bubble, not the underlying
geometry.

We can express this as a family of interpretation maps: for each domain $\mathcal{D}$,
there is a function $\iota_{\mathcal{D}}$ assigning to each abstract bubble
$B$ a domain-specific semantic object, and to each operator a domain-specific
transformation. The formal properties, admissibility conditions, and
inadmissibility pathologies transfer unchanged. We use the term "interpretation
map" rather than "functor" advisedly: while the construction has a
category-theoretic flavor, we do not here fix the morphism structure needed
for a full functorial account.

\subsection{Arithmetic and Computation}

The arithmetic interpretation has already been developed in
Section~\ref{sec:operators}. Here $\iota_{\text{arith}}(B_i) =$ an
expression scope, $U(B_i)$ measures remaining unevaluated subexpressions,
and Pop corresponds to evaluation of the innermost subexpression. The
programming interpretation is structurally identical: $\iota_{\text{prog}}(B_i)
=$ an execution context, corresponding to a frame on the call stack, and Pop
corresponds to a function return.

\subsection{Humor and Joke Structure}

A joke creates an initial expectation frame $B_1$, then builds narrative
momentum inside it. The punchline is not a simple Pop of $B_1$: it
retroactively reveals that the entire frame was nested inside a larger scope
$B_0$ whose existence was not previously signaled. The pop cascade
is $[B_0, B_1] \to [B_0'] \to [\;]$, but what produces the humor is the
sudden realization that $B_0$ existed all along. This is a conservative
reframe immediately followed by Pop: the topology of the joke is revised
in the moment of delivery, and then immediately collapsed. The relief and
pleasure of laughter corresponds to the sudden reduction of $L(\Sigma)$.

A callback joke adds a further layer: an earlier frame $B_{\text{old}}$ from
a previous joke, apparently closed, is revealed to still be open (or more
precisely, is reopened by the callback). This is closest to Meld: two
apparently independent frames are identified through a newly revealed
compatibility policy $\pi$.

\subsection{Harmonic Tension and Musical Cadence}

In tonal music, a harmonic progression opens expectation scopes that seek
resolution. The dominant seventh chord is perhaps the most familiar example:
it creates a strong expectation bubble whose resolution is the tonic. Extended
harmonic sequences build nested layers of tension at multiple temporal scales,
from the phrase to the section to the movement. A perfect authentic cadence
performs a Pop cascade across several of these layers simultaneously. A deceptive
cadence performs a conservative reframe: the expected Pop target (the tonic)
is replaced by an unexpected but grammatically admissible alternative,
preserving the ancestral harmonic logic while reassigning the immediate parent
of the resolution event. The emotional effect of a deceptive cadence, surprise
combined with recognition, corresponds precisely to the phenomenology of
conservative reframe.

\subsection{Proof Structure}

A mathematical proof opens the theorem's claim as a top-level scope. Lemmas
open subordinate scopes, each adding to the semantic load. The proof is
complete when every lemma scope has been closed by argument and the theorem
scope is closed by appeal to the assembled lemmas. The well-nestedness
constraint is not merely a stylistic convention of mathematical writing; it is
the structural condition that makes proofs valid. A proof that attempts to
invoke a lemma whose own proof is incomplete is formally inadmissible under
exactly the Pop admissibility condition. The experience of a proof clicking
into place is the experience of a large Pop cascade: the deepest auxiliary
results resolve first, their resolutions propagate outward, and the main
theorem's scope closes over a fully resolved landscape.


%% ── §6 ──────────────────────────────────────────────────────────────────────
\section{DAG Containment and the Geometry of Metaphor}
\label{sec:dag}

\subsection{The Limits of Tree Containment}

The tree-structured containment relation introduced in Section~\ref{sec:operators}
is sufficient for all cases examined so far. But it imposes a condition that
human cognition demonstrably violates: every bubble has at most one immediate
parent. Metaphor, analogy, and creative cross-domain reasoning require that a
single semantic object participate simultaneously in two or more containment
structures that are otherwise disjoint.

When we say \emph{argument is war}, we are not placing the domain of argument
inside the domain of war or vice versa. We are asserting that a semantic
region participates simultaneously in both containment structures, inheriting
properties from each. This multi-parent participation is not expressible in a
rooted tree. It requires a DAG.

\begin{definition}[DAG Containment Structure]
A \emph{DAG containment structure} is a pair $\Sigma = (B, \prec)$ where $B$
is a finite set of bubbles and $\prec$ is a strict partial order on $B$
without the single-parent restriction. The Hasse diagram $H(\Sigma)$ is a
directed acyclic graph.
\end{definition}

\subsection{Metaphorical Participation}

In a DAG containment structure, a bubble $B_i$ can have multiple immediate
parents $\pi_1, \pi_2, \ldots$ from disjoint containment chains. This
constitutes \emph{metaphorical participation}: $B_i$ simultaneously inherits
semantic structure from multiple domains.

This is not a meld operation. Meld closes two compatible sibling scopes into
one. Metaphorical participation maintains two parent scopes while the child
bubble remains open, inheriting from both without collapsing either. The
distinction is important: meld is post-resolution convergence, while
metaphorical participation is pre-resolution inheritance.

\subsection{Gluing Conditions and Metaphor Vitality}

Over a tree, the gluing condition (that local semantic sections be mutually
compatible) is nearly trivially satisfiable: each bubble has only one parent
to be compatible with. Over a DAG, gluing requires that a bubble's semantic
content be simultaneously compatible with all of its parent scopes.

This provides a precise criterion for metaphor vitality. A \emph{live
metaphor} is one in which the multi-parent gluing condition is satisfied but
nontrivially so: the bubble genuinely inherits from both parent domains, and
the overlap tension between them is semantically productive. A \emph{dead
metaphor} is one in which the multi-parent structure has collapsed into
conventional monadic participation: one parent has effectively become
vestigial, leaving the bubble with a single dominant parent and no productive
tension. The semantic curvature has flattened.

Metaphor quality thus corresponds to sustained admissible overlap without
forced collapse. The best metaphors maintain genuine multi-parent inheritance
over extended cognitive engagement. They resist the conventionalizing collapse
that would turn them into mere synonymy.

\subsection{Creative Cognition as Controlled DAG Expansion}

The expansive reframe operator $\Reframe^e_\phi$ takes on its fullest meaning
in the DAG setting. When a thinker introduces a new analogy or metaphor, they
are performing an expansive reframe on the current containment structure:
adding new ancestral dependencies that link previously disjoint domains. The
admissibility condition requires that the new structure remain acyclic and
that the gluing conditions for multi-parent bubbles remain satisfiable.

Creative insight at its deepest level may therefore consist in discovering
that a DAG expansion is admissible: that two previously disjoint containment
domains can be lawfully bridged without incoherence. The experience of
insight, simultaneous surprise and rightness, corresponds to the realization
that the new gluing condition is satisfied.


\subsection{Semantic Curvature and Attention Flow}

The semantic load functional $L(\Sigma)$ admits a geometric interpretation.
Unresolved semantic regions generate gradients in attentional space, bending
cognitive trajectories toward suspended dependencies awaiting closure. A highly
weighted unresolved bubble exerts stronger curvature, drawing working memory,
expectation, and interpretive effort toward itself even when not currently
foregrounded. This geometric interpretation is compatible with the physical
analogy suggested by the RSVP framework: just as regions of unresolved
field tension generate flow toward equilibrium, open semantic scopes generate
attentional flow toward resolution.

Under this interpretation, suspense corresponds to sustained unresolved
semantic curvature; curiosity corresponds to directed traversal toward
unresolved regions of the containment structure; distraction corresponds to
competing curvature wells drawing attention in incompatible directions; and
insight corresponds to a topological simplification that substantially reduces
global curvature complexity through a well-timed sequence of pops and
conservative reframes. The listener's experience of wanting to know what
happens next is thus not merely emotional anticipation. It is geometric
attraction within semantic space. Narrative tension, proof obligation, musical
expectation, and unfinished thought all manifest as unresolved curvature acting
on attentional flow. A practitioner of conversational hypnosis, by this
account, is sculpting a curvature landscape: opening regions of high semantic
load, maintaining their coherence, and timing the cascade of resolution for
maximum cognitive and affective effect.

This geometric reading also suggests a natural extension of the framework
toward quantitative cognitive modeling. The weight function $w_i =
f(d_i, s_i, r_i, c_i)$ already encodes the principal dimensions along which
semantic curvature varies. A full theory would supply empirical content for
$f$ — connecting stack depth, emotional salience, recency, and causal
centrality to measurable attentional and physiological signatures. That
empirical program lies beyond the scope of this paper but is a natural
downstream development of the formal framework.

%% ── §7 ──────────────────────────────────────────────────────────────────────
\section{Philosophical Consequences}

\subsection{Against Token-Linear Models of Cognition}

The framework developed here constitutes a direct argument against
token-linear models of symbolic cognition. If scope is primitive and not
derived, then any model that processes symbols sequentially without maintaining
persistent unresolved nested structure is cognitively inadequate in a
structural sense, not merely in a performance sense.

This is not an argument about computational power. A Turing-complete
token-linear system can simulate any containment structure. The argument is
about cognitive architecture: the natural grain of human symbolic cognition
is hierarchical containment, not sequential token succession, and models built
on the wrong grain will fail in characteristic ways.

\subsection{Premature Closure as Cognitive Failure Mode}

The inadmissibility condition on Pop identifies premature closure as the
most basic cognitive failure mode: resolving a scope before its dependencies
are discharged. This manifests as:
false certainty in reasoning (concluding before evidence is fully integrated),
narrative impatience (resolving a story's tension before its structure
warrants it),
therapeutic failure (accepting an explanation of distress before the
underlying frame has been examined), and
aesthetic crudeness (collapsing a work's productive ambiguity into
single-interpretation resolution before the audience has inhabited the
tension). In every case, the formal signature is the same: $\Pop(B_i)$ applied
while some descendant $B_j \prec B_i$ satisfies $U(B_j) > 0$.

\subsection{Intelligence as Disciplined Containment}

The deepest claim this framework supports is that intelligence is not
primarily prediction, inference, or pattern matching. It is the disciplined
management of nested unresolved semantic structure: knowing when to open a
scope, how to maintain it without premature closure, when to meld compatible
regions, and when a conservative or expansive reframe is warranted.

Ordinary push-pop systems treat Open and Pop as their primitive operations. While a Turing-complete stack machine can simulate Meld and Reframe computationally, it does not admit them as primitive semantic operators over containment topology: the operations must be encoded indirectly rather than expressed directly in the architecture's native grain. The cognitive claim is architectural, not computational. The capacity for Meld enables synthesis. The
capacity for conservative Reframe enables insight, therapy, and narrative
transformation. The capacity for expansive Reframe enables metaphor, analogy,
and creative cross-domain thinking. These capacities are not incremental
improvements on basic scope management. They are qualitatively distinct
operations that extend the operator basis in ways that standard stack
architectures cannot simulate at the cognitive level, even while simulating
them computationally.

\subsection{The Pop as Primitive of Meaning}

We close with the claim that motivates the entire framework. Meaning is not a
static relation between symbols and referents. It is the residue of an
irreversible act of semantic commitment: the closing of a scope that was
genuinely open. The value $5$ means something different from $(2 + 3)$ not
because they refer to different objects, but because one of them carries
provenance and the other carries only the compressed result of a history that
is no longer recoverable. Abstraction is not representation. It is the
quotient of a history of nested commitments by the relation of
resolution-equivalence.

This is why the joy of understanding feels the way it does. It is not the
passive recognition of a pattern. It is the active collapse of a structure
that was genuinely, productively, lawfully open.


%% ── Appendix ────────────────────────────────────────────────────────────────
\section{Toward Scope-Based Artificial Intelligence}

Most contemporary artificial intelligence systems are built around local
continuation dynamics over token sequences. While such systems can simulate
nested structure implicitly through learned statistical dependencies, they do
not generally maintain explicit persistent unresolved semantic scopes as
primitive architectural objects. The consequence is a characteristic
failure mode: premature closure. When local coherence pressures compete with
the demands of sustained unresolved structure, continuation-based systems tend
to collapse open scope too early, generating locally fluent but globally
incoherent output.

The framework developed here suggests a different architectural possibility:
intelligence as the lawful management of containment topology. Under such an
architecture, unresolved semantic regions would persist explicitly across
time as first-class objects; admissibility conditions would regulate collapse
operations, preventing premature resolution; and reframing would operate
directly on semantic inheritance structure rather than emerging indirectly from
statistical continuation weights.

The distinction is potentially foundational. Human cognition appears capable
of sustaining genuinely unresolved structure for extended durations while
dynamically reorganizing containment topology through conservative and
expansive reframing. Creativity, insight, metaphor, and deep narrative
understanding may depend less on predictive accuracy than on the controlled
maintenance and transformation of unresolved semantic geometry. A system that
can open a scope, hold it lawfully open under attentional pressure, meld it
with compatible sibling scopes, and reframe the containment hierarchy when
new structure becomes available would exhibit qualitatively different cognitive
behavior from a system whose architecture provides no native representation of
unresolved semantic load. This paper has developed the formal vocabulary for
specifying what such a system would need. What remains is the engineering.

The framework is already diagnostically useful for characterizing the failure
modes of existing systems. Hallucination corresponds to inadmissible expansive
reframe: new ancestral dependencies are introduced without satisfying gluing
conditions, producing locally fluent but globally incoherent output. Premature
summarization is illegal Pop: a scope is collapsed before its interior
dependencies are discharged. Sycophancy is forced Meld: two frames are
identified as compatible under social pressure rather than genuine
$\pi$-compatibility. Context-window degradation is scope-density overload:
when $D(\Omega)$ exceeds the system's representational capacity, containment
relations collapse indiscriminately. Mode collapse is semantic curvature
minimization: the system gravitates toward low-$L(\Sigma)$ configurations by
prematurely resolving all open scopes, eliminating the productive tension on
which creative and critical cognition depends. Incoherent long-range reasoning
is failure to preserve unresolved ancestry relations across time: the system
loses track of open obligations as the scope stack grows. Each failure mode
has a precise formal signature under the present framework, which suggests
that the framework could serve not merely as a design specification but as a
diagnostic instrument for evaluating existing architectures.

The present framework is compatible with dynamical accounts of meaning and
intentionality that treat causation as constraint propagation through evolving
constraint landscapes rather than linear efficient causation \cite{Juarrero1999}.
The inadmissibility of cyclic containment relations bears partial resemblance to
productive and pathological forms of recursive self-reference explored in
studies of strange loops \cite{Hofstadter1979}. The DAG treatment of metaphor
developed here overlaps partially with conceptual metaphor theory and
conceptual blending, though the present framework emphasizes containment
topology and admissible gluing rather than linguistic mapping alone
\cite{Lakoff1980,Fauconnier2002}. The containment interpretation shares
structural affinities with lambda calculus in its treatment of scope, binding,
reduction, and nested substitution \cite{Barendregt1984}, and with the
recursion-scheme tradition in functional programming \cite{Meijer1991}.
The categorical vocabulary of gluing, compositionality, and morphism structure
throughout the paper resonates with the foundations of category theory
\cite{MacLane1998}.


\subsection{Scope Closure and Temporal Directionality}

The asymmetry between open and closed scopes introduces an intrinsic temporal
orientation into the containment structure. An open bubble contains multiple
admissible future trajectories, while a collapsed bubble represents a reduced
equivalence class in which those trajectories have been quotient-compressed
into a resolved semantic residue. This gives the framework an internal arrow
of semantic time. The distinction between future and past is not imposed
externally upon the system but emerges from the irreversible reduction of
admissible possibility under lawful Pop operations. A resolved scope cannot be
reopened without either reconstructing its generative history or introducing a
new expansive reframe. The semantic directionality of cognition therefore
parallels thermodynamic irreversibility: both involve the progressive reduction
of accessible configuration structure under lawful evolution operators. The
analogy is structural rather than physical, but it suggests that the
irreversibility explored throughout this paper is not a special feature of
symbolic systems but an instance of a more general principle governing
constraint-governed dynamical processes.

\subsection{Local and Global Coherence}

The gluing condition introduced in connection with DAG containment admits an
important distinction between local and global semantic coherence. A
collection of bubbles may be pairwise compatible while still failing to admit
a globally coherent containment structure. Local compatibility does not
guarantee the existence of a globally consistent section over the entire
containment manifold. This distinction appears throughout cognition. Conspiracy
systems often maintain strong local coherence between neighboring explanatory
frames while failing globally under broader integration. Fragmented narratives
may preserve intelligibility scene by scene while collapsing under large-scale
interpretive traversal. Formally, many cognitive pathologies arise not from
local incoherence but from failures of global integrability: each adjacent
pair of scopes agrees, but no globally admissible configuration exists that
satisfies all constraints simultaneously. This is precisely the condition that
sheaf theory identifies as the failure of local sections to extend to a global
section, and it provides a precise mathematical characterization of a wide
class of reasoning failures that are otherwise difficult to describe.

\subsection{Semantic Inertia}

Not all containment structures reconfigure with equal ease. Certain bubbles
exhibit strong resistance to admissible reframing even when alternative
containment relations are available. We refer to this resistance as
\emph{semantic inertia}. Semantic inertia may arise from emotional salience,
repeated reinforcement, causal centrality, or dense integration with
surrounding scopes. Highly inertial bubbles resist both conservative and
expansive reframing because modification would propagate destabilization across
large portions of the containment graph. This provides a structural
interpretation of ideological rigidity, habit-formation, and persistent
self-concepts. Stability is not merely a psychological preference but a
topological property of deeply integrated semantic regions. An inertial bubble
functions as a load-bearing node: it is not merely heavy in the sense of high
$w_i U(B_i)$, but structurally central in the sense that its containment
relations are shared by many other bubbles. Reframing it requires
simultaneously satisfying compatibility conditions across all its dependent
scopes, which may be collectively unsatisfiable even when individually
admissible alternatives exist.

\subsection{Cascade Thresholds and Semantic Phase Transitions}

The semantic load functional $L(\Sigma)$ does not merely measure unresolved
structure statically. Beyond certain critical values, local operations may
propagate nonlocally through the containment graph, producing large-scale
reorganization. We therefore introduce a critical load parameter $L_c$ such
that a containment structure enters a \emph{cascade-sensitive regime} whenever
\[
  L(\Sigma) > L_c.
\]

\begin{definition}[Cascade Threshold]
A containment structure is \emph{cascade-sensitive} when $L(\Sigma) > L_c$
for some context-dependent critical load parameter $L_c$, beyond which
admissible local operations propagate nonlocally through the containment
graph.
\end{definition}

In a cascade-sensitive regime, a single admissible Pop, Meld, or Reframe may
trigger widespread restructuring across the containment topology. The collapse
of one heavily weighted bubble can suddenly reduce the unresolved status of
multiple dependent scopes, producing a rapid outward chain of semantic
closure. Humor, emotional breakthrough, sudden mathematical insight, and
certain forms of hypnotic release all exhibit this phenomenology: a prolonged
period of unresolved tension followed by abrupt large-scale reorganization.

The critical threshold $L_c$ is context-dependent and varies with the
stability of the surrounding containment structure. Highly coherent semantic
manifolds can sustain larger unresolved loads before cascading, while weakly
glued structures may collapse under comparatively small perturbations.
Semantic stability therefore depends not merely on total unresolved load but
on the geometry through which that load is distributed.

\subsection{Semantic Resonance and Recurrent Scope Activation}

Not all suspended bubbles remain equally active through time. Certain scopes
recur repeatedly across otherwise unrelated containment paths, reappearing as
persistent attractors in cognition, narrative, or emotional interpretation. A
\emph{resonant bubble} $B_r$ is characterized by repeated reactivation across
distinct regions of $\Sigma$:
\[
  B_r \rightsquigarrow \{ \Sigma^{(1)}, \Sigma^{(2)}, \ldots \}.
\]
The bubble's unresolved structure propagates into multiple future contexts,
modulating interpretation even after partial closure events have occurred. In
narrative, resonance appears as thematic recurrence or symbolic motif. In
trauma, a partially unresolved frame repeatedly reorganizes unrelated present
experience around itself. In mathematics, an unresolved conceptual tension may
quietly structure years of investigation before eventual resolution. In
creative work, resonance often precedes explicit insight: a scope continues
reappearing because the containment system has not yet discovered the
admissible reframe capable of integrating it coherently. Resonance acts as a
persistence amplifier: a highly resonant bubble maintains causal influence not
merely because it remains unresolved, but because the surrounding semantic
manifold repeatedly reconstructs paths that intersect it.

\subsection{Scope Density and Cognitive Compression}

Containment structures differ not only in total semantic load but in the
density with which unresolved dependencies are packed into local regions of
the graph. The \emph{scope density} of a region $\Omega \subseteq B$ is
\[
  D(\Omega) = \frac{\displaystyle\sum_{B_i \in \Omega} U(B_i)}{|\Omega|}.
\]
High-density regions contain many unresolved dependencies compressed into a
small topological neighborhood. Such regions exhibit heightened cognitive
pressure, rapid attentional cycling, and increased susceptibility to cascade
behavior. This helps explain why certain symbolic systems feel cognitively
compressed: a highly abstract proof, a densely metaphorical poem, or a
multilayered hypnotic induction may carry extremely high local scope density
within small surface structure. Cognitive fluency therefore depends not solely
on total load $L(\Sigma)$ but on the geometric distribution of unresolved
structure across the containment manifold.

\section{Constraint Propagation and Semantic Causality}

The framework implies a conception of causality different from linear
push-pull models of symbolic processing. Within a containment structure,
unresolved scopes constrain the admissible future evolution of the system long
before any explicit collapse occurs. A bubble with high unresolved load
restricts which future Pops, Melds, and Reframes remain admissible. The
structure therefore evolves not merely through local symbol transitions but
through globally distributed constraint propagation. Future semantic
trajectories are shaped by the topology of currently unresolved dependencies.

This interpretation aligns naturally with dynamical approaches to cognition
in which causation operates through evolving constraint landscapes rather than
isolated efficient interactions \cite{Juarrero1999}. The unresolved scope acts
less like a stored object and more like a curvature condition imposed on future
semantic motion. Meaningful cognition is therefore prospective as well as
retrospective. Open scopes organize the future before they collapse into the
past.

\section{Toward a Vector-Valued Load Theory}

The semantic load functional $L(\Sigma) = \sum_i w_i U(B_i)$ has served as a
useful scalar summary of unresolved semantic tension throughout this paper.
But the weight decomposition $w_i = f(d_i, s_i, r_i, c_i)$ already signals
that semantic load is not a single quantity. Different dimensions of weight
can dissociate: a mathematically unresolved problem may carry high structural
load and low affective salience. A trauma trigger may carry enormous affective
curvature despite shallow logical structure. A conspiracy system may maintain
high affective coherence simultaneously with low causal coherence. The scalar
$L(\Sigma)$ compresses these dissociations into a single value, which is
adequate for the foundational treatment but insufficient for a full cognitive
theory.

A natural extension replaces the scalar load with a vector-valued functional:
\[
  \mathbf{L} : \Sigma \to \mathbb{R}^n,
  \quad
  \mathbf{L}(\Sigma) = (\lambda_s L_s,\; \lambda_a L_a,\; \lambda_c L_c,\; \lambda_p L_p),
\]
where $L_s$ measures structural dependency load, $L_a$ measures affective
salience load, $L_c$ measures causal centrality load, and $L_p$ measures
phenomenological tension, each weighted by context-dependent coupling
constants $\lambda_s, \lambda_a, \lambda_c, \lambda_p$. Under this
decomposition the containment manifold admits different curvature along
different load dimensions. A system can be geometrically flat in the causal
dimension while exhibiting steep affective gradients, producing the
dissociation characteristic of intellectualized trauma or dispassionate
mathematical obsession. The cascade threshold condition $L(\Sigma) > L_c$
generalizes naturally to a threshold surface in $\mathbb{R}^n$, allowing
different cognitive systems to exhibit phase transitions along different
projections of the load manifold. This extension remains a research program
rather than a completed formalism, but the foundational definitions of this
paper are consistent with it and in several places anticipate it.

\subsection{Homotopy Classes of Derivation Histories}

The collapse quotient $B/{\sim_\rho}$ introduced in Section~\ref{sec:operators}
raises the question of how much of a derivation history survives compression.
The weakest interpretation is purely extensional: two derivations are
equivalent iff they collapse to the same final residue, so that $(2+3)$,
$(1+4)$, and $(10/2)$ become indistinguishable after resolution. This is
computationally natural but phenomenologically inaccurate. Different
derivational paths leave behind structural traces that affect subsequent
cognition. A theorem discovered geometrically behaves differently from the
same theorem discovered algebraically, even after the formal proof is
identical. A traumatic insight reached gradually behaves differently from one
reached catastrophically, even if the explicit conclusion is verbally
indistinguishable afterward.

The stronger interpretation treats the collapse residue as preserving a
compressed homotopy class of the derivation history:
\[
  B/{\sim_\rho^\alpha}
\]
where $\alpha$ is a resolution parameter controlling the coarseness of the
equivalence. At coarse $\alpha$, all derivational paths to a given result
are identified. At finer $\alpha$, derivations partition into semantic
equivalence classes distinguished by their topological structure, emotional
trajectory, modal character, or reframing depth. This stratification permits
the framework to distinguish fact equivalence from experiential equivalence
— a distinction that appears essential for any adequate theory of identity,
expertise, creativity, or style. The full development of homotopy-stratified
collapse quotients is a natural direction for subsequent formal work.

\section{Recursive Identity and Narrative Continuity}

Human identity may be interpreted within this framework as a persistent,
partially unresolved containment structure maintained across time through
continuous cycles of Pop, Meld, and Reframe. A self is not a static object
but an evolving topology of unresolved commitments, inherited narratives,
anticipated futures, and recursively reorganized memories.

Under this interpretation, autobiographical continuity emerges from the
preservation of admissible ancestry relations across repeated reframing
operations. A person changes not by replacing one fixed identity with another
but by continuously redistributing containment relations over a partially
stable historical graph. Conservative reframes preserve identity through
reinterpretation: $\prec^* = \prec'^*$. The same historical dependencies
remain present while their local organization changes. Expansive reframes
incorporate previously external domains into the identity graph:
$\prec^* \subseteq \prec'^*$. New ancestral structures become integrated into
the self-model.

This provides a natural account of both psychological continuity and personal
transformation. Stability arises from preserved ancestry; growth arises from
topological reorganization. Identity is therefore neither immutable essence
nor arbitrary reconstruction. It is recursive containment maintained through
lawful semantic evolution. Pathological rigidity corresponds to a failure to
admit any admissible Reframe. Pathological dissolution corresponds to a
failure to maintain any stable ancestry relation across time. Health, under
this account, is the sustained capacity for both.


\section{The Self as Deep Attractor Bubble}

The conversation between Karl Friston and Susan Blackmore on predictive
processing and the free energy principle \cite{Friston2023} converges on several claims that are
structurally precise translations of the containment framework developed in
this paper. The convergence is not merely analogical. The two frameworks are
approaching the same geometry from different entry points: predictive processing
from the perspective of Bayesian brain theory and active inference, and the
present framework from the perspective of recursive semantic containment and
deferred closure.

Friston's central claim about the self is that it is not a controller sitting
at the apex of a strict hierarchy but a deep, enduring hypothesis embedded
within a heterarchical structure. In his account, the deeper a representation
sits within the generative model, the more slowly it changes and the more
temporally extended its influence. The self is the deepest such representation
precisely because it must remain stable across the widest range of incoming
evidence. It is not at the top because it commands; it is at the center
because it persists.

This maps directly onto the containment framework. A bubble with very high
temporal persistence and very low $U(B)$ turnover rate functions as an
attractor in semantic space: it organizes surrounding scopes without being
quickly resolved by them. The self-bubble $B_{\text{self}}$ is therefore
characterized not by high unresolved load in the ordinary sense but by
extremely high semantic inertia. Its containment relations are shared by an
enormous portion of the surrounding graph. Reframing it propagates
destabilization throughout the entire structure, which is precisely why
identity change is experienced as difficult, disorienting, and effortful.

The heterarchical structure Friston proposes — self at the center rather than
the top, with predictions propagating centrifugally outward and evidence
flowing centripetally inward — is structurally equivalent to replacing the
rooted-tree model of self with a DAG model in which $B_{\text{self}}$ is not
the root but a highly connected interior node. Its influence is not top-down
command but multi-parent containment inheritance: it shapes the admissibility
conditions for a large portion of the surrounding semantic graph without
directly controlling any single node.

Friston introduces the concept of \emph{mortal computation} (a phrase due
originally to Geoffrey Hinton) to characterize computation that is
substrate-dependent in a constitutive rather than incidental way. On a von
Neumann architecture, memory and processing are separated, allowing the
software to be copied, distributed, and run on arbitrary hardware. The
computation is therefore abstract: detached from any particular physical
trajectory. On a mortal computational substrate — neuromorphic, memristive,
or photonic — the computation is the physical process. It cannot be copied
without being ended, because the substrate and the computation are the same
thing.

This distinction maps onto the irreversibility framework of this paper. A pop
operation in the Spherepop calculus is irreversible precisely because the
derivation history is not separately stored. The collapse of $(2+3)$ into $5$
is not the deletion of a record; it is the ending of a process. The history
existed as a trajectory, not as a log. Under a mortal computational substrate,
this is not a limitation but the constitutive structure of the process. The
trajectory \emph{is} the computation. There is no copy. This provides a
formal account of why von Neumann architectures cannot, in Friston's sense,
truly self-evidence: self-evidencing requires that the organism's current
model state be shaped by its own prior physical history in an unbroken
trajectory of constraint propagation. A system that can be copied, reset, or
run on an arbitrary substrate has no such trajectory. Its semantic bubbles
carry no mortal provenance.

\section{Suffering, Precision, and Recursive Self-Containment}

Friston's account of suffering is structurally precise: to suffer, an organism
must include within its generative model a hypothesis that \emph{I am a self
who can suffer}. Suffering is not bare negative affect; it is negative affect
recursively bound to a self-model with future-continuity constraints. Without
that recursive binding, an organism may exhibit instability, error-signaling,
or avoidance behavior without experiencing anything. A thermostat registers
deviation from setpoint without suffering, because there is no persistent
hypothesis that \emph{I am the kind of thing that should be at setpoint and
whose failure to be so matters}.

This can be formalized within the scope framework. Let $\Pi_{\text{self}}$
denote the precision weighting — in the active inference sense, the inverse
variance — assigned to the self-hypothesis bubble $B_{\text{self}}$. Let
$U(B_{\text{future}})$ denote the unresolved semantic load of the
future-directed scopes currently nested inside $B_{\text{self}}$. Then
suffering in its structural sense can be approximated as:
\[
  \mathcal{S} \;\sim\; \Pi_{\text{self}} \cdot U(B_{\text{future}}).
\]
High precision weighting on the self-hypothesis means that deviations from
expected self-consistent outcomes register strongly. High unresolved forward
load means that many anticipated scopes remain open. The product of these two
quantities produces the characteristic phenomenology of suffering: not mere
uncertainty, but self-indexed unresolvedness. Ordinary uncertainty does not
hurt unless it is recursively bound to a persistent self-model that cannot
tolerate the openness.

This formulation resolves a puzzle in pure predictive processing accounts.
Expected surprise (entropy) is technically bad in the free energy framework,
but organisms clearly tolerate uncertainty in many contexts without suffering.
The resolution is that entropy without self-binding is merely load. It becomes
suffering when $\Pi_{\text{self}}$ is high enough to amplify $U(B_{\text{future}})$
into a cascade-sensitive regime. Below the threshold $L_c$, unresolved
forward load is manageable. Above it, the cascade of unresolved
self-indexed scopes produces the recursive amplification characteristic of
anxiety, grief, and existential distress.

Grief, as Friston and Blackmore discuss it, is a specific instance of this
structure. A major loss renders the self-model no longer fit for purpose: the
forward-directed scopes that were organized around a particular world structure
(domestic life, relational patterns, spatial familiarity) are suddenly
inadmissible, because the world they anticipated no longer exists. The
relearning required is not merely the accumulation of new facts but the
replacement of large portions of the containment graph while $B_{\text{self}}$
must somehow remain stable enough to sustain the process. It is exhausting
because it is semantically expensive: an enormous number of previously closed
scopes must be reopened, re-evaluated, and re-collapsed under the new
constraints, all while the self-bubble that organizes them is itself under
reconstruction.

\section{Meditation, Top-Layer Relaxation, and the Dissolution of Recursive Agency}

Blackmore describes the meditative process in terms that map precisely onto
the containment framework: taking the top layers off, reducing the precision
weighting on high-level hypotheses, and discovering what remains when the
recursive self-model is no longer actively maintained. Friston interprets this
as a reduction in the precision weighting of deep priors:
$\Pi_d \to 0$ for the deepest layers of the generative hierarchy.

In scope-dynamic terms, meditative dissolution involves the progressive
relaxation of semantic inertia at the highest levels of the containment
structure. The self-bubble $B_{\text{self}}$ does not pop in the ordinary sense
— that would be psychotic dissolution, not meditative insight — but its
precision weighting decreases, reducing its gravitational pull on surrounding
scopes. The forward-directed load $U(B_{\text{future}})$ decreases not because
the scopes are forcibly closed but because the self-hypothesis that amplified
their urgency is no longer strongly weighted.

The result, as both Friston and Blackmore describe, is not emptiness but a
shift in the geometry of attentional flow. When $\Pi_{\text{self}}$ decreases,
the semantic curvature generated by $B_{\text{self}}$ flattens. The attentional
gradients that ordinarily bend cognition toward self-relevant unresolved scopes
relax. What remains is local scope navigation without the global self-indexed
amplification. Perception continues; cognition continues; but the recursive
binding that converts unresolved load into suffering is temporarily suspended.

Conversely, pathological states correspond to excessive $\Pi_{\text{self}}$:
the self-bubble becomes so inertially dominant that its containment relations
propagate into every region of the semantic manifold, rendering all
unresolved scopes self-relevant and therefore unbearable. Anxiety disorders,
certain psychotic presentations, and existential crises all share this
structure: the cascade threshold is effectively lowered to zero because
$\Pi_{\text{self}}$ amplifies every open scope into a potential crisis.

The framework thus provides a unified account of ordinary selfhood, suffering,
grief, meditative dissolution, and pathological self-amplification as different
parameter regimes of the same underlying structure: recursive containment under
varying precision weighting. This is not a replacement for predictive
processing theory but a geometric restatement of its central claims in the
vocabulary of scope dynamics. The two frameworks are complementary. Predictive
processing specifies the computational mechanism — Bayesian inference over
hierarchical generative models. The scope-dynamic framework specifies the
topological structure within which that inference operates — nested unresolved
regions under admissible collapse operators. Together they describe the same
system from different angles of approach.

\appendix

\section{Formal Definitions Collected}

For reference, the principal definitions of the paper are collected here in
order of introduction.

A \emph{semantic bubble} $B = (C, E, U)$ consists of a contextual binding set
$C$, an expectation structure $E$, and a non-negative unresolved load scalar
$U(B)$.

A \emph{containment structure} $\Sigma = (B, \prec)$ is a finite set of
bubbles equipped with a strict partial order. In the tree case every bubble
has at most one immediate parent. In the DAG case this restriction is dropped.

The \emph{scope stack} $\Sigma_t = [B_1, \ldots, B_n]$ linearizes the active
containment path with $B_n$ the currently attended scope.

The \emph{semantic load functional} is $L(\Sigma) = \sum_i w_i U(B_i)$ with
$w_i = f(d_i, s_i, r_i, c_i)$.

The \emph{resolution operator} $\rho : B \times \Sigma \rightharpoonup B'$
is partial, defined iff all descendants of $B$ are closed.

\emph{Well-nestedness} at position $i$ requires $U(B_j) = 0$ for all $B_j
\prec B_i$.

The five primitive operators are $\Open$, $\Pop$, $\Meld_\pi$,
$\Reframe^c_\phi$, $\Reframe^e_\phi$, with admissibility conditions as
stated in Section~\ref{sec:operators}.

\emph{Conservative reframe} satisfies $\prec^* = \prec'^*$.
\emph{Expansive reframe} satisfies $\prec^* \subseteq \prec'^*$.
Both require $\prec'$ to be acyclic.

\section{Worked Examples}

\subsection*{Arithmetic}

$( 3 + ( 4 \times ( 2 + 1 ) ) )$ reduces as
\[
  (3 + (4 \times 3)) \;\to\; (3 + 12) \;\to\; 15.
\]
Each arrow is an admissible Pop. The innermost scope $(2+1)$ pops first
because its interior is empty; its parent $(4 \times 3)$ becomes eligible
only after that pop; and so on outward.

\subsection*{Narrative Induction}

Open $B_1$ (frame story), interrupt to open $B_2$ (embedded story), interrupt
to open $B_3$ (innermost anecdote). Resolve: $B_3$ pops, modifying $B_2'$;
$B_2'$ pops, modifying $B_1'$; $B_1'$ pops, emptying the stack. Total
semantic load decreases monotonically through the resolution cascade.

\subsection*{Conservative Reframe: Therapeutic}

A client holds a childhood experience $B_{\text{trauma}}$ as the enclosing
frame for their adult identity $B_{\text{self}}$. After therapeutic
reframing, $B_{\text{trauma}}$ becomes a contained episode within a broader
developmental narrative $B_{\text{history}}$. Node set unchanged; ancestral
graph unchanged; covering relations redistributed. The event remains real;
its containment role changes.

\subsection*{Expansive Reframe: Metaphor}

The domain of argument $B_{\text{arg}}$ acquires an additional parent scope
$B_{\text{war}}$ through the metaphor \emph{argument is war}. The transitive
closure $\prec^*$ is extended. $B_{\text{arg}}$ now inherits from both its
original parent structures and the new war-domain parent. Gluing requires
that the inherited properties from $B_{\text{war}}$ be simultaneously
compatible with $B_{\text{arg}}$'s existing parent constraints.


%% ── Bibliography ────────────────────────────────────────────────────────────
\begin{thebibliography}{99}

\bibitem{Barendregt1984}
Henk Barendregt.
\textit{The Lambda Calculus: Its Syntax and Semantics}.
North-Holland, 1984.

\bibitem{Baez2020}
John Baez and Bob Coecke, editors.
\textit{Applied Category Theory 2019}.
Electronic Notes in Theoretical Computer Science, 2020.

\bibitem{Fauconnier2002}
Gilles Fauconnier and Mark Turner.
\textit{The Way We Think: Conceptual Blending and the Mind's Hidden
Complexities}.
Basic Books, 2002.

\bibitem{Hofstadter1979}
Douglas Hofstadter.
\textit{G\"{o}del, Escher, Bach: An Eternal Golden Braid}.
Basic Books, 1979.

\bibitem{Juarrero1999}
Alicia Juarrero.
\textit{Dynamics in Action: Intentional Behavior as a Complex System}.
MIT Press, 1999.

\bibitem{Lakoff1980}
George Lakoff and Mark Johnson.
\textit{Metaphors We Live By}.
University of Chicago Press, 1980.

\bibitem{Meijer1991}
Erik Meijer, Maarten Fokkinga, and Ross Paterson.
Functional programming with bananas, lenses, envelopes and barbed wire.
In \textit{FPCA '91: Functional Programming Languages and Computer
Architecture}, pages 124--144. Springer, 1991.

\bibitem{Friston2023}
Karl Friston and Susan Blackmore.
Why do models suffer? Consciousness, self-evidencing, and mortal computation.
Transcript of conversation, 2023.

\bibitem{MacLane1998}
Saunders Mac Lane.
\textit{Categories for the Working Mathematician}.
Springer, 2nd edition, 1998.

\bibitem{Needham1997}
Tristan Needham.
\textit{Visual Complex Analysis}.
Oxford University Press, 1997.

\end{thebibliography}

\end{document}
